\documentclass[11pt]{article}
\usepackage[T1]{fontenc}
\usepackage[utf8]{inputenc}
\usepackage{lmodern}
\usepackage[margin=1in]{geometry}
\usepackage{amsmath,amssymb,amsthm,mathtools}
\usepackage{microtype,etoolbox,xcolor,url,needspace}
\usepackage[titles]{tocloft}
\setlength{\cftbeforesecskip}{2.8pt}
\usepackage[colorlinks=true,linkcolor=blue!50!black,citecolor=green!35!black,urlcolor=blue!60!black]{hyperref}
\usepackage[capitalize,nameinlink]{cleveref}
\hypersetup{pdftitle={Linear-Time Tape Transposition Requires Storage Codes with Rates Approaching One},pdfauthor={Hans Brende and GPT-6 Astra}}
\theoremstyle{plain}
\newtheorem{theorem}{Theorem}[section]
\newtheorem{lemma}[theorem]{Lemma}
\newtheorem{proposition}[theorem]{Proposition}
\newtheorem{corollary}[theorem]{Corollary}
\theoremstyle{definition}
\newtheorem{definition}[theorem]{Definition}
\newtheorem{remark}[theorem]{Remark}
\newenvironment{informalproof}[1][Written proof]{\begin{proof}[#1]}{\end{proof}}
\newenvironment{outline}[1][Outline of the formal proof]{\begin{proof}[#1]}{\end{proof}}
\newcommand{\notformalized}{\par\noindent{\normalfont\small\textsf{Proof status.} Written proof.}\par\smallskip\noindent\ignorespaces}
\DeclareUrlCommand\leanname{\def\UrlFont{\ttfamily\upshape}}
\renewcommand*\UrlBreaks{\do\_\do\.\do\/\do\-}
\newcommand{\lean}[1]{\leanname{#1}}
\newcommand{\leanfile}[1]{\leanname{#1}}
\newcommand{\leanlist}[1]{\begingroup\def\leansep{}\renewcommand*{\do}[1]{\leansep\leanname{##1}\def\leansep{, }}\docsvlist{#1}\endgroup}
\newcommand{\formal}[2]{\par\noindent{\normalfont\small\raggedright\textsf{Lean.} \leanlist{#2} in \leanfile{#1}.\par}\smallskip\noindent\ignorespaces}
\newcommand{\N}{\mathbb{N}}
\newcommand{\Z}{\mathbb{Z}}
\newcommand{\R}{\mathbb{R}}
\newcommand{\eps}{\varepsilon}
\newcommand{\card}[1]{\lvert #1\rvert}
\newcommand{\abs}[1]{\lvert #1\rvert}
\newcommand{\tpos}{\operatorname{tpos}}
\DeclareMathOperator{\dist}{dist}

\setlength{\emergencystretch}{2em}
\setcounter{tocdepth}{1}
\BeforeBeginEnvironment{theorem}{\Needspace{8\baselineskip}}
\BeforeBeginEnvironment{proposition}{\Needspace{8\baselineskip}}
\BeforeBeginEnvironment{corollary}{\Needspace{8\baselineskip}}
\clubpenalty=3500
\widowpenalty=3500
\title{Linear-Time Tape Transposition Requires Storage Codes with Rates Approaching One}
\author{Hans Brende \and GPT-6 Astra (OpenAI)}
\date{November 2026}
\begin{document}
\maketitle
\begin{abstract}
Suppose a deterministic multitape machine transposes square matrices in
linear time at arbitrarily large dimensions $d$, with entries of
$\max(1,\lceil\log_2 d\rceil)$ bits and head moves that may depend on the input. For every sufficiently large incidence-degree
cap $\Delta$ we obtain a family of binary storage codes whose minimum
distance between paired input and recovery ports tends to infinity. Its
limiting lower rate is at least
\[
  1-O_{M,c}\!\left(
       \left(\frac{\log_2\log_2\log_2\Delta}{\log_2\log_2\Delta}\right)^{1/(4k^2+1)}
     \right)
\]
for $k$ work tapes and time coefficient $c$. The transfer also holds on
jointly growing dimensions and entry widths when the rate of the successful
input set tends to one. For codes on the original bit positions, recovered at their actual
append times on a simulator, the degree-dependent loss has exponent
$1/(8k^2+2)$. On any nonempty successful set $G$ of $L$-bit inputs, the
fraction of inputs whose fibre rate falls more than this allowance plus
$u\ge0$ below $\log_2|G|/L$ is at most $2^{-uL}$. For these codes, bounded-overlap patches face a further
restriction. The sum of the information their permitted helpers supply about
individual separated targets is at most twice the overlap times the total
correlation of the input.
All numbered results are proved in Lean. An upper bound on the attainable
storage rates at the required scale would exclude linear-time transposition.

\end{abstract}
\begin{quote}\small
\textbf{Contributions.} Hans Brende chose the problem and directed the research. GPT-6 Astra developed the arguments, wrote the manuscript and formalized the proofs in Lean.
\end{quote}
\tableofcontents
\clearpage
\section{Introduction}\label{sec:introduction}

The entries of a matrix's first column are a row apart on the input tape
and consecutive in the output. A transposer may store coded information
as it works, and its head movements may depend on what it has read. If
we fix those movements, we can trace the possible dependencies of an
output bit. Fixing the movements also restricts the input. The question
is how much information survives the restriction.

For a fixed tape trajectory, a depth-reduction theorem supplies a small
set of steps whose values can be revealed so that the remaining
dependency paths have bounded length \cite[Corollary~O]{PaperFourteen}.
On the inputs sharing both the trajectory and the revealed values, each
output can then be recovered from a bounded set of input coordinates.
The recovered coordinate and its input copy can be far apart in the
incidence graph of that relation. Such separation is compatible with
binary storage codes of rate above one half
\cite[Corollary~A$'$]{PaperThirteen}. A linear-time transposer would
require rates arbitrarily close to one.

\Cref{thm:main} gives the degree dependence of this requirement. For
every sufficiently large degree cap it produces a family with rates
bounded below while its minimum paired distance tends to infinity. The
coordinates are bits of packets. Grouping bits this way gives our
strongest rate estimate. A coding upper bound with a larger deficit would
rule out the assumed machine. \Cref{cor:gap} makes the comparison precise,
including the success probability when the machine is correct only on
some inputs.

A second question concerns the individual bit positions of a physical
machine. \Cref{thm:simulator} groups source transitions into batches but
executes each head move and each append on an ordinary bit machine.
\Cref{thm:native-tail} uses its actual append times. The packet graph connects every bit of a permitted helper packet to
every bit of its target packet. The native graph follows each physical
append backwards through the simulator. On a set $G$ of correctly
transposed inputs, the native theorem bounds the input mass in fibres
whose rate falls appreciably below $\log_2|G|/L$. Each fibre keeps nearly
all of its trace's original bit positions. The native exponent is half
the packet exponent.

Finally, \cref{sec:patches} asks how much a local set of helpers can tell
us about the bit it is meant to recover. The answer accounts separately
for a bit's own bias and the information supplied by its helpers. A
fixed advantage over chance on a positive fraction of separated targets
requires increasingly large patch overlap as the code rate approaches
one. The outstanding question is whether a linear-time tape computation
would have to provide such an advantage at smaller overlap.

\paragraph{Relation to earlier work.}
The finite-control simulation of Hartmanis and Stearns evaluates
consecutive tape operations from a bounded neighbourhood
\cite[Theorem~2]{HartmanisStearns}. Hopcroft, Paul and Valiant synchronize
tape-block crossings \cite[\S2]{HPV1977}. Williams describes the
computation graph of an aligned machine using a constant number of
movement bits per block \cite[\S\S2--3]{Williams}. Here the sequence of batch descriptors determines the in-batch append
flags and when the terminal drain starts. The common output length
then fixes the number of bits appended during the drain. Earlier in this series, an external-memory
simulation fixed an event word at the scale of entire rows
\cite[Theorem~EM]{PaperEleven}. Our batch size is fixed before the
matrix dimension grows. Conditional entropy as a measure of imperfect
local recovery appears in earlier work on recoverable systems
\cite[\S IV]{ElishcoBarg}. Our prediction bound uses the separated-port
argument of \cite[Lemma~4.1]{PaperFive} and the soft-recovery inequality
for nested elimination \cite[Theorem~S]{SoftRecovery}.

\paragraph{Formal proofs.}
The Lean declarations named with each numbered result are its proofs.
The English outlines explain the constructions, estimates and choices of
parameters. \Cref{sec:verification} gives the verification procedure. The declaration index identifies the corresponding source modules.

\section{Tape computations and recovery relations}\label{sec:model}

\paragraph{The machine and its task.}
We use the deterministic multitape model of the first paper in the series
\cite{PaperOne}. The machine has finite control, a two-way read-only input
tape, $k\ge0$ two-way work tapes over a fixed finite alphabet, and an
append-only binary output. All heads start at zero and move by at most one
cell per transition. A transition writes at the scanned work cells and
appends at most one bit. Reading beyond the input gives an end marker.
The \emph{full motion trace} records the head directions and whether a bit
was appended at each transition, together with termination. The next control state, written symbols and appended bit belong to the
step value \cite[\S2]{PaperEleven}.

A $d\times d$ matrix with $b\ge1$ bits per entry occupies $L=d^2b$ bits in
row-major order. For $0\le i,j<d$ and $0\le t<b$, the input bit at address
$q=(id+j)b+t$ must be appended at output position
$\tpos(q)=(jd+i)b+t$. Fix a machine $M$ and a time coefficient $c$.
Write $G_{d,b}$ for the inputs on which $M$ halts within $c(L+1)$ steps
with this output. At the canonical width $b(d)=\max(1,\lceil\log_2 d\rceil)$,
statement \textnormal{(T)} asks for a single $M,c$ with
$G_{d,b(d)}=\{0,1\}^L$ for every positive $d$.

\paragraph{A code and its paired distances.}
Let $V$ be a finite nonempty coordinate set and $\mathcal A$ a finite
alphabet of size at least two. A relation $N$ assigns to each $q\in V$
a set $N(q)\subseteq V$ of permitted helpers. A nonempty set
$C\subseteq\mathcal A^V$ is a \emph{storage code} for $N$ if two words of
$C$ agreeing on $N(q)$ agree at $q$, for every $q$. Its rate is
$\log_{|\mathcal A|}|C|/|V|$. No linearity is assumed.

The incidence graph has an input port $L_v$ and a recovery port $R_v$
for each $v$. Join $L_v$ to $R_q$ when $v\in N(q)$. The \emph{paired
distance} of $q$ is the undirected graph distance from $L_q$ to $R_q$,
with value $\infty$ for disconnected ports. For example, take two
coordinates that store the same bit and let each use its partner as its
only helper. The two incidence edges are disjoint. Both paired distances
are infinite, although each coordinate is recovered exactly. This graph
measurement concerns the permitted observations at the two ports. The example has rate $1/2$ \cite[\S3]{PaperFive}.

A binary family $(V_i,N_i,C_i)$ is \emph{admissible at cap $\Delta$} if
every port has degree at most $\Delta$, $|V_i|\to\infty$, and for each
fixed integer radius $r$ every paired distance exceeds $r$ for all
sufficiently large $i$. Let $R(\Delta)$ be the supremum, over these
families, of their limiting upper rates. Our constructions give lower
bounds on the limiting \emph{lower} rate. For large $\Delta$, write
\[
 p=4k^2+1,\qquad
 f_j(\Delta)=\left(\frac{\log_2\log_2\log_2\Delta}
                         {\log_2\log_2\Delta}\right)^{1/j}\quad(j>0).
\]
The packet and native bounds below use $j=p$ and $j=2p$, respectively.

\paragraph{How a tape run gives a recovery relation.}
Freeze a motion trace. In the diagram of step dependencies, a step has
at most $k+1$ earlier predecessors, one for the previous control step
and one for the last visit to the scanned cell on each work tape. Missing predecessors
give initial values. The input head supplies the bit at its current
address, or the end marker. Choose a set $\Gamma$ of steps and reveal
their full values. This collection of observations is called a
\emph{genie}. Start at the append of output bit $\tpos(q)$ and follow
the predecessor edges backwards, requiring every vertex after the root to lie outside $\Gamma$. The root
itself may be revealed.
The input addresses encountered form the \emph{cone} of that append.
Once the trace and the revealed values have been fixed, agreement on
these input bits forces agreement at the append, by induction on time.
On correctly transposed inputs, it therefore forces agreement at $q$.
This is the recovery relation used at physical bit scale. The same
recursion with packet inputs and packet records will be used in
\cref{sec:proof} \cite[\S4]{PaperEleven}.

For $k$ work walks with $T\le C(n+1)$ steps, $C\ge1$, $n\ge1$, a cut
budget $\varepsilon>0$ and records of width $\beta\ge1$, the bounded-stack theorem gives a cut with
$\beta|\Gamma|\le\varepsilon n$. Each permitted path then has at most
$D+1$ vertices, where $D=\operatorname{genieDepth}(k,\beta,C,\varepsilon)$
is independent of $n$ \cite[Corollary~O(a)]{PaperFourteen}. Here $n$
counts input units. A unit is one packet in \cref{sec:proof} and one
physical bit in \cref{sec:native}. In both applications the revealed
record values are charged in bits.

\section{The rate forced by a packet computation}\label{sec:proof}

\begin{theorem}[Theorem A, quantitative packet transfer]\label{thm:main}
\formal{TapeStorageRates/Main.lean}{TapeStorageRates.rate_witnesses_of_growing_widths,TapeStorageRates.rate_witnesses_of_near_full,TapeStorageRates.rate_witnesses_of_full}
Fix a deterministic machine $M$ with $k$ work tapes and a time coefficient $c>0$.
Suppose that for every $n\ge1$ and $\eta>0$ there are $d,b\ge n$ for which
\[
  G_{d,b}\ne\varnothing,\qquad
  \log_2|G_{d,b}|\ge(1-\eta)d^2b .
\]
There is a constant $K=K(M,c)>0$ such that for every sufficiently large
integer $\Delta$ some admissible binary family at cap $\Delta$ satisfies
\[
  \liminf_{i\to\infty}\frac{\log_2|C_i|}{|V_i|}
     \ \ge\ 1-K f_{4k^2+1}(\Delta).
\]
Consequently $R(\Delta)\ge1-K f_{4k^2+1}(\Delta)$ for all sufficiently
large $\Delta$, and $R(\Delta)\to1$.
\end{theorem}
\begin{outline}
Fix $s$ bits per packet. On a sufficiently large successful instance,
conditioning on the incomplete ends of entries, the packet schedule and
the revealed steps leaves a code with a deficit of $O_{M,c}(1)$ bits
per original packet. The trimming estimate below retains at least half
the packets at the chosen cap. Expanding those packets into bits gives
a rate loss of $O_{M,c}(1/s)$. The required degree cap $\Delta_s$ satisfies
$\log_2\log_2\Delta_s=O_{M,c}(s^p\log_2(s+1))$. Inverting it gives
the claimed rate. The following two subsections establish the finite
estimates and make the order of the choices explicit.
\end{outline}

Both thresholds $d,b\ge n$ apply to the same success set. In particular,
the hypothesis holds if $M$ succeeds on $2^{L-o(L)}$ inputs along a sequence
of canonical dimensions tending to infinity. A machine satisfying (T)
succeeds on all $2^L$ inputs at those dimensions. The family in the
conclusion may depend on $\Delta$. Once the cap is fixed, its length
and its minimum paired distance grow together.

\subsection{From successful inputs to a common packet schedule}

Fix a packet size $s\ge1$. Use the hypothesis of \cref{thm:main} with $\eta=1/(2s)$.
We can choose $d$ as large as needed with $b\ge s$ and a nonempty success
set $G=G_{d,b}$ of deficit $h=L-\log_2|G|\le L/(2s)$. Write $b=as+r$
with $a\ge1$ and $0\le r<s$. Each matrix entry consists of $a$ complete
$s$-bit packets followed by $r$ tail bits. There are $n=d^2a$ complete
packets. The packet at $(i,j,\ell)$ is demanded at $(j,i,\ell)$ for $0\le\ell<a$. Since $L=sn+rd^2\le2sn$, we have $h\le n$.

Partition $G$ by the entire string of $rd^2$ tail bits. Some fixed
assignment $z$ leaves a set $F_z$ of successful packet arrays with
\[
  |F_z|\ge |G|\,2^{-rd^2}=2^{sn-h}.
\]
Thus fixing the tails preserves the bound $h$ on the absolute deficit,
with $sn$ as the new ambient bit count. Throughout this computation the
tail of an entry is the fixed string specified by its address in $z$.

The auxiliary input tape uses one cell per packet. The last cell of each
entry holds its $s$ payload bits followed by its fixed tail, so its length
in the original input is $s+r$. All other in-range cells have length $s$.
Each cell outside the input represents $s$ end markers, and each work cell
holds $s$ ordinary work symbols. Starting from any offset in a cell, $s$
source transitions can visit only that cell and its two neighbours.

One batch evaluates these transitions in eleven packet ticks. The input
and work heads move one cell left and read the left, centre and right
cells in succession. A stationary tick evaluates the batch on these
windows. The heads traverse the windows in reverse to write back the
three work cells, return to the old centres, and move by at most one cell
to the new centres. Other work writes preserve the scanned symbols. The final stationary tick may append one packet.
At the start of a batch the pending output is shorter than its next
cell. The batch adds at most $s$ raw output bits, so it completes at
most one cell. On completion, emit its first $s$ bits, pass over its
tail, and keep the suffix for the next cell. Every successful run emits
exactly the $n$ payload packets in transposed order. The public cell
lengths and tail values belong to this auxiliary computation.
\Cref{sec:simulation} will return to ordinary bit tapes.

Use $R=\lfloor\lfloor c(L+1)\rfloor/s\rfloor+1$ batches on every input,
keeping the source configuration fixed after halt. For
$J=2\lceil c\rceil+1$, the bounds $R\le J(n+1)$ and $T=11R\le11J(n+1)$
are uniform in $s,d,b$. Within each batch the varying part of the
schedule consists of the $k+1$ displacements to the new centres and the
packet-append flag. These give at most $(2\cdot3^{k+1})^R$ schedule
words. Hence for an integer $A_0=A_0(k,c)$ some fixed schedule has a
class $W$ satisfying
\[
  \log_2|W|\ \ge\ sn-h-A_0 n\ \ge\ sn-(A_0+1)n.
\]
The compiler further records the finite control, offsets, windows,
pending output, work writes and any output packet. They fit in
$\beta_s=O_M(s)$ bits per tick. On the fixed schedule the input address,
cell type, tick within the batch, and count of previously appended
packets are public. The record at a tick is determined by its input
packet, the preceding control record, and the records from the last
visits to the scanned work cells. Initial values and exterior input
symbols are fixed. This locality and the output at each designated
append tick are proved by \lean{TapeStorageRates.PaidClasses.compiler}.

\subsection{Extracting and separating a code}

Call $W$ a \emph{frozen packet class} on $n=d^2a$ packets if its nonempty
set of packet arrays has common unit walks for the input head, the $k$
work heads and the output position, and its tick records obey the
locality rule above. It also has a designated tick before the deadline
for each demanded packet. At that tick the output position is the
packet's transposed address and the record determines the packet on
$W$. In the compiled class this is its append tick.

\begin{proposition}[Retention at a chosen degree threshold]\label{prop:packet-retention}
\formal{TapeStorageRates/PacketFamily.lean}{TapeStorageRates.packet_code_at_threshold}
Fix integers $k,A\ge0$, $s,\beta\ge1$ and a time coefficient $c_0\ge1$.
Let $D=\operatorname{genieDepth}(k,\beta,c_0,1)$ be the depth bound
of the bounded-stack walk theorem, and put
$B=\sum_{j=0}^{D}(k+1)^j$. Fix a degree threshold $\theta\ge0$, a
radius $u\ge0$, both integers, and a real $\alpha$ with
$0<\alpha<1-B/(\theta+1)$.
For all sufficiently large $d$, every $a\ge1$ and every frozen class $W$
on $n=d^2a$ packets with $\beta$-bit records, at most $c_0(n+1)$ ticks
and $\log_2|W|\ge sn-An$, there is an integer $m$ with $\alpha n\le m\le n$ and a binary
storage code on $sm$ coordinates.
Its degree is at most $s\max(B,\theta)$, every paired distance exceeds
$u$, and its rate is at least $1-(A+1)/(\alpha s)$.
\end{proposition}

\begin{outline}
Point each tick back to its control and work predecessors. The walk
theorem supplies a set $\Gamma$ of ticks with $\beta|\Gamma|\le n$.
On a dependency path starting at an output tick and avoiding $\Gamma$
after its first vertex, the depth bound is $D$.
Fix a most frequent joint value for the records at $\Gamma$. This
costs at most $n$ bits in log cardinality. For a source packet $q$, let
$N(q)$ contain the input addresses reached by these paths from its
designated output tick. Two arrays in the resulting fibre that agree on
$N(q)$ have the same output record there. Indeed the revealed records
already agree, and the locality rule gives agreement at the remaining
ticks by induction in time. Thus $N(q)$ recovers $q$. This application
uses the whole-step walk cut of \cite[Corollary~O(a)]{PaperFourteen} and
the common-schedule recovery argument of \cite[proof of Corollary~6.4]{PaperEight}.

Each tick has at most $k+1$ predecessors. Consequently $|N(q)|\le B$,
and the sum of all input-port degrees is at most $Bn$. Remove both ports
of each coordinate whose input degree exceeds $\theta$. There are at
most $Bn/(\theta+1)$ such coordinates. In the induced graph the maximum
degree is at most $\max(B,\theta)$. The paired-path packing theorem bounds
the number of remaining matched pairs at distance at most $u$ by
$O_{k,c_0,D,\theta,u}(n/\log_2 d)$, uniformly in $a$. At bounded degree, a positive density of short paths would yield many
vertex-disjoint ones. Following their dependency edges forwards and
backwards through the walks would transport the corresponding input
coordinates to their transposed positions. The row-entropy transport
bound limits the number that can be served in $O(n)$ ticks. This is
where the transposed addresses enter the estimate. The formal packing
bound uses the walks, the cut and the positions of the transposed demands
\cite[Theorem~3.3 and Remark~3.9]{PaperFive}. At block scale one, the Lean estimate
\lean{MergeDepth.skel_paired_packing} underlying
\cite[Proposition~4.11(a)]{PaperEleven} also permits $k=0$. The strict inequality defining $\alpha$
leaves room to remove these short pairs and retain $m\ge\alpha n$.
Every surviving path was already a path before deletion.

Fix the omitted packets to common values when projecting the code through
the two deletions. The total cost is at most $s(n-m)$ bits, one $s$-bit
payment for each omitted coordinate. Its former occurrences in the
recovery rules are now constants \cite[Corollary~3.8]{PaperFive}. The
projected code has log cardinality at least $sm-(A+1)n$.
Replace every surviving packet by its $s$ bits and each incidence by
all $s^2$ incidences between the corresponding copies. The bits of the
helper packets determine each demanded bit. Expansion preserves code
cardinality, multiplies the degree bound by $s$, and sends every paired
path back to a packet walk of the same length when its bit indices are
forgotten \cite[Lemma~8.1]{PaperThirteen}. Dividing the cardinality
bound by $sm$ gives the rate in the statement.
\end{outline}

\begin{outline}[Completion of the proof of Theorem A]
Apply \cref{prop:packet-retention} with $c_0=11J$, $\beta=\beta_s$ and
$A=A_0+1$ from the schedule count. Taking $\theta=2B$ permits
$\alpha=1/2$, since $B/(2B+1)<1/2$. Thus the cap is $\Delta_s=2sB$
and the rate loss is at most $2(A+1)/s$. At cap $10sB$, taking $\theta=10B$ permits
a retained fraction of at least $9/10$.
For each successive radius $u$, take a qualifying $d,b$ large enough
for the proposition and with $d\ge2(u+1)$. Their lengths tend to infinity, and
the $u$th code has every paired distance greater than $u$.

It remains to fit $\Delta_s$ below a prescribed $\Delta$.
Write $x=4\beta_s c_0=O_{M,c}(s)$. In the walk cut the full depth is
$D=K_0^{P_0}\lceil x\rceil$, with its radix $K_0$ and digit count $P_0$
taken at cut fraction $1/x$ \cite[Definition~7.1]{PaperFourteen}.
For fixed $k$ the parameter bounds give
$\log_2 K_0=O_k(x)$ and
$P_0=O_k(x^{4k^2}\log_2(x+1))$.
Since $B\le(k+2)^{D+1}$, these inequalities yield
\[
   \log_2\log_2\Delta_s
      =O_{M,c}\!\left(s^{\,4k^2+1}\log_2(s+1)\right).
\]
With $y=\log_2\log_2\Delta$ and $p=4k^2+1$, choose $s$ as the
integer part of a sufficiently small fixed multiple of
$(y/\log_2 y)^{1/p}$. For all sufficiently large $\Delta$ this gives
$\Delta_s\le\Delta$ and $1/s=O_{M,c}(f_p(\Delta))$.
The family at that $s$ proves the limiting lower rate bound.
The supremum $R(\Delta)$ is at least this family's limiting upper rate,
and every binary code has rate at most one. This proves the convergence
asserted in \cref{thm:main}.
\end{outline}


\section{A degree gap sufficient for exponential failure}\label{sec:consequences}

Theorem~A gives a test for an impossibility proof. At a degree cap where
all admissible families have a rate gap $\rho$, a putative transposer
must pay it from its deficit allowance $K f_p(\Delta)$. The machine
constant $K$ is fixed before $\Delta$ is chosen.

\begin{corollary}[A sufficient gap at a sequence of caps]\label{cor:gap}
\formal{TapeStorageRates/Main.lean}{TapeStorageRates.exponentially_small_at_widths_of_sparse_gap}
Fix a tape cap $k\ge0$ and put $p=4k^2+1$. Suppose there are integer
caps $\Delta_j\to\infty$ and real numbers $\rho_j$ such that
$\rho_j/f_p(\Delta_j)\to\infty$. At each $\Delta_j$, assume that
every admissible family has limiting lower rate at most $1-\rho_j$.
For each fixed machine $M$ with at most $k$ work tapes and each $c>0$,
there are $\gamma>0$ and $n_0$ such that
\[
   \frac{|G_{d,b}|}{2^{d^2b}}
      \le 2^{-\gamma d^2b}
      \qquad(d\ge n_0,\ b\ge n_0).
\]
In particular the same exponential upper bound holds at every
sufficiently large canonical dimension, and (T) fails for such machines.
\end{corollary}

\begin{outline}
If the joint near-full success hypothesis of \cref{thm:main} held,
each sufficiently large cap would admit a family of limiting lower rate
at least $1-K f_p(\Delta)$. A machine with fewer than $k$ tapes also
satisfies this bound, since increasing the tape count increases the
deficit allowance at large caps. Choose $j$ with $\Delta_j$ large
enough for the lower bound and $\rho_j>K f_p(\Delta_j)$. The required
family violates the assumed upper bound at that same cap.
Negating the near-full success hypothesis now gives some $\gamma>0$ and
$n_0$ such that, throughout $d,b\ge n_0$, each nonempty success set has
$\log_2|G_{d,b}|\le(1-\gamma)d^2b$. Empty success sets satisfy the
probability bound as well. Finally $b(d)\to\infty$ with $d$.
\end{outline}

For example, suppose $R(\Delta)\le1-\rho(\Delta)$ at every sufficiently
large $\Delta$, with $\rho(\Delta)\ge0$. Writing
$y=\log_2\log_2\Delta$, a gap $\rho(\Delta)\ge c_1/y^a$ for a constant
$c_1>0$ and $0<a<1/p$ suffices for the fixed tape cap $k$.
If $\limsup_{\Delta\to\infty}\rho(\Delta)y^a=\infty$ for every $a>0$,
the exclusion applies to every fixed finite tape count. Indeed
$y^{1/(2p)}f_p(\Delta)\to0$. A gap bounded below by a positive
constant times $1/\log_2 y$ meets this all-tapes condition.
The elementary gap $1/(\Delta+1)$ obtained by ordering coordinates
has ratio to $f_p(\Delta)$ tending to zero
\cite[open problems]{PaperThirteen}. The slowly decreasing gap criterion
for oblivious machines appears after Corollary~O in \cite{PaperFourteen}.
The power comparison and the elementary-gap comparison are checked in
\lean{TapeStorageRates.ratio_tendsto_of_power_gap} and
\lean{TapeStorageRates.elementary_gap_ratio_tendsto_zero}.
The coding upper bound at the scale required for \cref{cor:gap} remains open.

\section{A bit simulator with few motion traces}\label{sec:simulation}

We now keep the matrix in its original binary layout. The batch size
controls how often a simulator reveals an input-dependent decision in
its head motions. The purpose is to leave a large set of inputs on which
the physical cone at each append has the same meaning.

\begin{theorem}[Theorem S, native trace count]\label{thm:simulator}
\formal{TapeStorageRates/NativeCompiler.lean}{TapeStorageRates.Native.same_partial_function,TapeStorageRates.Native.source_support_fixed,TapeStorageRates.Native.source_support}
Fix a deterministic machine $M$ with $k$ work tapes and an integer $s\ge1$.
There is a machine $N_s$ with the same number of tapes and the same work alphabet that computes the same partial function. On a halting run of $t$ successful transitions of $M$, its running time is exactly
\[
 (10s+1)(\lfloor t/s\rfloor+1)+s.
\]
Its finite control can be encoded in $O_M(s)$ bits. On any set of halting inputs with $t\le U$ and common output length, its number of distinct full motion traces is at most
\[
 (1+H_k)^{\lfloor U/s\rfloor+1},\qquad H_k=2\cdot3^{k+1}.
\]
If the output lengths vary, multiply this bound by $s$.
\end{theorem}
\begin{outline}
At the start of a batch the physical head on tape $j$ sits at an anchor $p_j\in s\Z$. Its simulated head is at $p_j+r_j$, where $0\le r_j<s$. The simulated work tapes agree with the physical ones. The control holds the offsets, simulated state and a word of fewer than $s$ pending output bits. The physical output followed by this word is the simulated output so far. All anchors and offsets start at zero.

Move the heads left $s$ cells. In the next $3s$ ticks read the windows $[p_j-s,p_j+2s-1]$, moving right except on the final read. In one stationary tick a finite table simulates up to $s$ transitions on these windows and stops if it detects halt. Every visited cell and the final position are inside the scanned window, since the starting offset is in $[0,s-1]$. Each ending offset relative to the old anchor has a unique expression $se_j+r'_j$, where $e_j\in\{-1,0,1\}$ and $0\le r'_j<s$. The table also joins the new output bits to the pending word. Their total length is at most $2s-1$. If there are at least $s$ bits, schedule the first $s$ for emission.

Write the updated work windows back from right to left in $3s$ ticks, moving left except on the last write. Every head follows this sweep, including the input head. Return each head to its old anchor in $s$ rightward ticks. In the next $s$ ticks each takes direction $e_j$, reaching its new anchor. Hold the heads still for $s$ ticks and append the scheduled bits, or idle throughout if none were scheduled. The pending word is the unprinted suffix. These phases preserve the boundary conditions. Work writes outside the writeback phase preserve the scanned symbol.

In Lean, \lean{TapeStorageRates.Native.simulates_halt} proves the exact clock and equality of completed outputs. Each batch costs $10s+1$ ticks. If a batch discovers halt, drain the pending word during $s$ final stationary ticks, idling after its last bit, then halt. When $s$ divides $t$, the extra batch discovers halt before simulating any further transition. An infinite simulated run executes infinitely many batches. The windows, pending output, offsets, counters and finite state admit $O_M(s)$ bits of control by \lean{TapeStorageRates.Native.control_bits}. For fixed $M,s$ the update on this finite set is a transition table.

A batch's $k+1$ directions $e_j$ and its full-word emission flag determine all its native moves and append flags. Pad the sequence of batch descriptors, using a distinct symbol, to length $\lfloor U/s\rfloor+1$. Its unpadded length determines when to drain and halt. At a common output length $\ell$ the drain emits exactly $\ell\bmod s$ bits. At variable lengths this count takes at most $s$ values. This gives the claimed trace bounds.
\end{outline}

For inputs of length $L$ with $t\le c(L+1)$ and common output length, the limiting upper value of $L^{-1}\log_2|\operatorname{supp}(\mathrm{trace})|$ is at most $c\log_2(1+H_k)/s$. The limsup inequality is \lean{TapeStorageRates.Native.trace_support_rate_limsup}. Thus any prescribed positive trace rate selects one simulator before $L$ grows.

The time coefficient can also be made uniform in $s$ on all halting inputs. Simulate the first $2s$ successful transitions directly, stopping if $M$ halts earlier. In the batched phase, use the current physical head positions as grid origins and start with empty pending output. For $t<2s$ the total cost is $t$. Otherwise it is
\[
 3s+(10s+1)(\lfloor t/s\rfloor-1)\le 11t .
\]
The variant uses $O_M(s)$ bits of control. Its prefix adds $O_k(s)$ bits
to the trace description. The prefix flags specify how many bits were
already emitted there, so a fixed total output length still determines
the length of the final drain. The formal bounds are
\par{\small\raggedright
\lean{TapeStorageRates.Prefix.simulates_halt},
\lean{TapeStorageRates.Prefix.time_le_eleven},
\lean{TapeStorageRates.Prefix.control_bits}, and
\lean{TapeStorageRates.Prefix.source_support}.\par}

For this compiler the order $1/s$ can be attained, as proved in \lean{TapeStorageRates.Native.copy_trace_card}. Consider a copying machine that advances a dummy work head exactly on a $1$, and advances the input head after each copied bit. Independently choose each of the first $\lfloor L/s\rfloor$ aligned input blocks to be all zeros or all ones, fixing the suffix. In each batch the dummy head stays at its anchor or advances to the next anchor. These choices give $2^{\lfloor L/s\rfloor}$ different simulator traces. The example measures the compiler's trace count on a specified machine. Copying also admits an oblivious implementation.

\section{Typical codes on physical bit positions}\label{sec:native}

Run the simulator $N_s$ of \cref{thm:simulator} on a nonempty set
$G\subseteq G_{d,b}$. We will count inputs in $G$ with equal weight. The
construction has to keep track of which physical bit positions survive,
because the most frequently read bits can differ between traces.

\subsection{Counting the cells of a frozen trace}\label{sec:records}

For each attained full trace $\tau$, choose one set of revealed steps
$\Gamma_\tau$. Root the cone for input coordinate $q$ at the step that
appends output bit $\tpos(q)$. Delete both ports of each coordinate
whose input port has excessive degree. In the induced graph delete both
ports of every remaining coordinate whose paired distance is at most the
chosen radius. Denote the retained input positions by $U_\tau$ and set
$m_\tau=|U_\tau|$. The choices of steps and positions are common to the
whole trace. The existence of such cuts with bounded depth is proved in
\lean{TapeStorageRates.native_typical_partition} using the walk theorem
of \cite[Corollary~O(a)]{PaperFourteen}.

For a fixed $\tau$, record the values at $\Gamma_\tau$ and the input bits
outside $U_\tau$. Call this complete record $z$. Let $G_{\tau,z}$ be
the inputs with that trace and record, and $C_{\tau,z}$ their
projections onto $U_\tau$. Fixing the exterior makes the projection
injective. It also turns every omitted helper into a constant, so
$C_{\tau,z}$ is a code for the induced relation on the retained ports.
Each codeword extends to a member of $G_{\tau,z}$.
These are the storage and extension statements of
\lean{TapeStorageRates.NativeTypicalCut.storage} and
\lean{TapeStorageRates.NativeTypicalCut.mem_code_iff}. The full trace fixes the successive append flags, hence the step at
which each output position is written.

\begin{proposition}[Counting native fibres]\label{prop:moment}
\formal{TapeStorageRates/NativeTypical.lean}{TapeStorageRates.NativeTypicalPartition.deficit_moment,TapeStorageRates.NativeTypicalPartition.tail_of_joint_count}
Let $X$ be uniform on $G$ and use a native partition as above. For an
attained cell define its deficit
$D_{\tau,z}=m_\tau-\log_2|C_{\tau,z}|$. Write $D_X$ for the deficit of
the cell containing $X$. For any real budget $\ell_{\tau,z}$ on each
attained record, with $\ell_X$ the corresponding random budget,
\[
 \mathbb E\,2^{D_X-\ell_X}
   =\frac{1}{|G|}\sum_\tau\sum_{z:\,G_{\tau,z}\ne\varnothing}
                2^{m_\tau-\ell_{\tau,z}} .
\]
If at most $2^h$ distinct pairs of trace and revealed step values occur
on $G$, then for every $t\ge0$,
\[
 \Pr\!\left[D_X>L-\log_2|G|+h+t\right]\le 2^{-t}.
\]
\end{proposition}
\begin{outline}
The $|G_{\tau,z}|$ inputs in a cell all contribute the same power of two.
Since $|C_{\tau,z}|=|G_{\tau,z}|$, this cardinality cancels when we sum
over the cell, giving its term in the identity.

For the second display, use the constant budget
$\ell=L-\log_2|G|+h+t$. A fixed trace and fixed revealed word allow at
most $2^{L-m_\tau}$ exterior assignments. Their combined contribution
to the unnormalized sum is at most $2^{L-\ell}$, independently of
$m_\tau$. Summing over the attained pairs gives
$\mathbb E\,2^{D_X-\ell}\le2^{-t}$. On the event $D_X>\ell$ its
integrand exceeds one. The same argument with the individual budgets
bounds $\Pr[D_X>\ell_X]$ when the set of possible exterior assignments
is known more precisely.
\end{outline}

\subsection{A tail bound at fixed degree}\label{sec:native-rate}

\begin{theorem}[Theorem N, typical native rate]\label{thm:native-tail}
\formal{TapeStorageRates/NativeTypicalRates.lean}{TapeStorageRates.typical_native_rate_tail,TapeStorageRates.typical_native_rates_full}
Fix $M$ with $k$ work tapes and $c\ge0$. Put $p=4k^2+1$ and
$Z_\Delta=2^{\sqrt{\log_2\Delta}}$. There are constants $K,A>0$ such
that at every sufficiently large integer cap $\Delta$ one can fix
an integer $s\ge1$ with
\[
       f_{2p}(\Delta)\le 1/s\le K f_{2p}(\Delta).
\]
For every fixed integer radius $r\ge0$, all sufficiently large $d$,
every $b\ge1$ and every nonempty $G\subseteq G_{d,b}$, the bit simulator
$N_s$ has a partition of $G$ with the following properties.

Projecting an attained cell as in \cref{sec:records} gives a storage code
on the retained input positions of its trace. Every
input and recovery port has degree at most $\Delta$, every surviving paired distance
exceeds $r$, and
\[
               m_\tau\ge (1-Z_\Delta/\Delta)L .
\]
The simulator halts within $(11c+1)(L+1)$ steps on $G$. For $X$ uniform
on $G$, let $R_X=\log_2|C_{\tau,z}|/m_\tau$ for its cell.
On this same partition, for every $u\ge0$,
\[
 \Pr\!\left[R_X<\frac{\log_2|G|}{L}-A f_{2p}(\Delta)-u\right]
      \le 2^{-uL}.
\]
\end{theorem}
\begin{outline}
Let $\beta_s$ be the number of bits sufficient for a scale-one step
record of $N_s$. The simulator keeps its tape alphabet, and its control
uses $O_M(s)$ bits, so $\beta_s=O_M(s)$. At fixed $s$ the startup
cost fits in $L+1$ steps for all sufficiently large $d$, uniformly in
$b\ge1$. The exact clock and trace count in \cref{thm:simulator} then
give the displayed time bound and at most $2^{aL/s}$ traces on $G$, for
a constant $a=a(M,c)$. All inputs in $G$ have the same output length.

On each trace spend at most $L/s$ bits on revealed steps. In the cut
theorem the parameter controlling the depth is
$x=4\beta_s(11c+1)s=O_{M,c}(s^2)$
\cite[Definition~7.1 and Corollary~9.3(a)]{PaperFourteen}. If $D_s$ is
the resulting depth parameter, put $B_s=\sum_{j=0}^{D_s}(k+1)^j$.
This bounds the number of permitted input helpers at each output.
The degree calculation of \cref{sec:proof} now gives
$\log_2\log_2(2s^2 B_s)=O_{M,c}(s^{2p}\log_2(s+1))$.
Choose the constant in the inversion so that this double logarithm is
at most $\tfrac12\log_2\log_2\Delta$. Then
$\max(s,B_s)\le Z_\Delta$ at the stated scale for $s$, as proved in
\lean{TapeStorageRates.native_batch_size_for_cap}. This extra factor
of $s$ in the cut parameter accounts for the change from $f_p$ to
$f_{2p}$.

There are at most $B_sL$ incidences. Removing input coordinates of degree
greater than $\Delta$ therefore removes at most $B_sL/(\Delta+1)$
coordinates. In the induced graph, the short-pair estimate bounds the
second deletion by $J L/\log_2 d$ for a constant $J$ depending on the
fixed parameters and radius, uniformly over the traces and the widths
$b\ge1$ \cite[Proposition~4.11(a), accompanying Lean theorem]{PaperEleven}.
Since $B_s/(\Delta+1)<Z_\Delta/\Delta$, both deletions fit in the
claimed fraction at all sufficiently large $d$. Deleting coordinates
can only increase the remaining paired distances.

Put $\lambda=Z_\Delta/\Delta$. For large caps $Z_\Delta^2\le\Delta$,
so $\lambda\le1/s$. The chosen traces and genie words jointly cost at
most $(a+1)L/s$ bits. Apply \cref{prop:moment} with $t=uL$. Outside
an event of probability at most $2^{-uL}$ we obtain
\[
 \log_2 |C_{\tau,z}|
 \ge\log_2|G|-\bigl(\lambda+(a+1)/s+u\bigr)L .
\]
The code is nonempty, so its rate on $m_\tau\le L$ coordinates is at
least $\log_2|C_{\tau,z}|/L$. The loss
$\lambda+(a+1)/s$ is at most $(a+2)K f_{2p}(\Delta)$. This proves the
tail with a single partition for every confidence level.
\end{outline}

For a full successful cube $\log_2|G|/L=1$. Taking $u=1/s$ gives rate
$1-O_{M,c}(1/s)$ except on input mass at most $2^{-L/s}$. For a smaller
success set its own rate enters with coefficient one. The retained
positions and recovery relation may vary from trace to trace. The code can also vary with the revealed word and the exterior values.

\subsection{A sufficient native coding bound}\label{sec:native-test}

The native half-rate hypothesis concerns the full cone relation of a
specified machine at a specified depth. It asks for an upper bound
$2^{(1/2+\delta)L}$ on every binary storage code on each eligible cone
relation for all sufficiently large $d$, for every fixed $\delta>0$.
The step width $\operatorname{blockWidth}(N_s,1)$ is the number of bits
charged for revealing one complete physical record.

\begin{corollary}[The native frozen-trace criterion]\label{cor:native}
\formal{TapeStorageRates/TraceTransfer.lean}{TapeStorageRates.native_frozen_trace_exclusion,TapeStorageRates.native_tests_exclude_linear}
For each $M$ and $c>0$ one can choose a single aligned simulator $N_s$ and
time coefficient $c'>0$ with the following implication. Suppose the
native half-rate hypothesis of \cite[Corollary L(b)]{PaperFourteen}
holds at canonical widths for $N_s$ at scale one, time coefficient
$c'$, and depth parameter $\operatorname{genieDepth}
(k,\operatorname{blockWidth}(N_s,1),\max(c',1),1/8)$. Then $M$
succeeds on the full cube at only finitely many canonical dimensions.
If this hypothesis holds on the chosen simulator for each source machine
and coefficient, \textnormal{(T)} is false.
\end{corollary}
\begin{outline}
An eligible run halts by the specified time bound. The hypothesis
ranges over cuts with the prescribed
remaining depth and all choices of times before halt at the prescribed
output-counter positions. In particular it includes the actual append
steps. Choose $s$ so that the complete trace support costs at most $L/4$
bits at large successful dimensions, using
\lean{TapeStorageRates.aligned_low_trace_dimensions}, which supplies
$c'=11c+1$. A largest trace
contains at least $2^{3L/4}$ inputs. Revealing at most $L/8$ bits on
that trace leaves a common word with at least $2^{5L/8}$ inputs. These
inputs form a code on the untrimmed append-time cones by the same
predecessor induction as in \cref{sec:model}. This exceeds the hypothesized
bound for every $0<\delta<1/8$. The formal criterion applies the skeleton-entropy lower bound of
Corollary~L(b). Skeleton entropy is at most the logarithm of the full
trace support. The separate witness theorem
\lean{TapeStorageRates.native_five_eighths} constructs the native code
and \lean{TapeStorageRates.NativeFibreWitness.outputTime_eq_appendTime}
identifies its roots.
\end{outline}

The same checked criterion allows any fixed $0<\eta<1/2$. Choose the
simulator so that its full trace support costs at most $\eta L$ bits
at large successful dimensions. In the depth parameter, replace the
cut fraction $1/8$ by $(1/2-\eta)/2$. The common trace and genie fibre
then has rate at least $1-\eta-(1/2-\eta)/2>1/2$.

\section{What can overlapping patches predict?}\label{sec:patches}

A bit that is zero with probability \(0.9\) can be predicted with \(90\%\)
accuracy without reading any helper. The information supplied by an allowed
helper set \(B\) is \(I(X_q;X_B)=H(X_q)-H(X_q\mid X_B)\). In the prediction credit
\[
  1-H(X_q\mid X_B)=[1-H(X_q)]+I(X_q;X_B),
\]
the first term is due to the marginal bias and the second measures the
gain from observing the helpers. The geometry of the relation bounds
the sum of these gains, even when different
patches use the same bits. Write \(X_S=(X_v)_{v\in S}\) and use entropy
\(H\) in bits. The Shannon entropy in the Lean declarations is measured in
nats, so its value in bits is obtained by division by \(\ln 2\).

Let \(V\) be a finite set of coordinates. Each member of a finite indexed family of \emph{patches} \(P_j\subseteq V\)
has a set \(A_j\subseteq P_j\) of scored targets. For each \(q\in A_j\),
choose the helpers \(B_{jq}\subseteq P_j\) available to its predictor.
A target scored in two patches contributes two occurrences \((j,q)\).
Form the bipartite graph on the two copies of \(A_j\) by joining \(L_v\) to
\(R_q\) when \(v\in B_{jq}\cap A_j\). A patch satisfies the
\emph{separation condition} if \(L_q\) and \(R_q\) lie in different
components of this graph for every \(q\in A_j\).

An ambient distance bound gives a sufficient separation test. Take a
relation \(N\). Suppose \(B_{jq}\subseteq N(q)\cap P_j\),
\(\card{A_j}\le w\) with \(w\ge1\), and
\(\dist(L_q,R_q)>2w-1\) for every scored occurrence. A connection inside
the patch would contain a simple path on at most \(2\card{A_j}\) ports.
The same path in the ambient graph would have length at most \(2w-1\).
Thus the separation condition holds. The patch can contain arbitrarily
many unscored helpers.

\begin{theorem}[Theorem P, information supplied by overlapping patches]
\label{thm:patch-credit}
\formal{TapeStorageRates/PatchResidual.lean}{TapeStorageRates.weighted_patch_capacity,TapeStorageRates.patch_capacity_of_long_targets}
Let \(h\) be a polymatroid rank on \(V\), that is, a normalized, monotone,
submodular set function. Write \(h(q\mid B)=h(B\cup\{q\})-h(B)\).
Give each coordinate a capacity \(c_q\ge h(\{q\})\) and each patch a
weight \(\lambda_j\ge0\). Suppose the separation condition holds, and
\(\sum_{j:v\in P_j}\lambda_j\le r\) for every coordinate \(v\). Then
\[
 \sum_j\lambda_j\sum_{q\in A_j}\bigl(c_q-h(q\mid B_{jq})\bigr)
 \le r\left[\left(\sum_{v\in V}c_v-h(V)\right)
              +\left(\sum_{v\in V}h(\{v\})-h(V)\right)\right].
\]
For a finite random vector \(X\), write
\(\mathrm{TC}(X)=\sum_v H(X_v)-H(X_V)\) for its total correlation.
Taking \(h(S)=H(X_S)\) and \(c_v=H(X_v)\) gives
\[
 \sum_j\lambda_j\sum_{q\in A_j} I(X_q;X_{B_{jq}})
 \le 2r\,\mathrm{TC}(X).
\]
\end{theorem}
\begin{outline}
Fix one patch and write \(P=P_j\), \(A=A_j\), \(D=P\setminus A\).
Condition on \(D\). All remaining helpers of a target \(q\) have their
left ports in the component of \(R_q\). Order the components and first
supply the targets whose left component precedes their right component.
Then process the other targets in order of their left components. For
each such target its right component came earlier, so all of its helpers
are available. A supplied target costs at most \(h(q\mid D)\) in the
chain rule, and a target reached with its helpers costs at most
\(h(q\mid B_{jq})\). Reversing the order exchanges the two roles for
each target. Adding the bounds gives
\[
 2\bigl(h(P)-h(D)\bigr)
 \le \sum_{q\in A} h(q\mid D)+\sum_{q\in A}h(q\mid B_{jq}).
\]
The formal proof groups the targets by their left components to obtain
a single elimination piece in the sense of~\cite{PaperSeven}. The
common-observations corollary of the soft-recovery memo gives this
inequality~\cite[Corollary~3.2]{SoftRecovery}. The two component orders explain the
same factor two as in the exact reconstruction argument
of~\cite[Lemma~4.1]{PaperFive}.

Use \(h(q\mid D)\le h(\{q\})\) and bound \(h(D)\) both by
\(\sum_{v\in D} h(\{v\})\) and by \(\sum_{v\in D} c_v\).
The contribution of this patch is at most the sum of the two deficits
on the right side of the theorem, evaluated on \(P\).
To sum either deficit, let \(a_v\) stand for \(c_v\) or \(h(\{v\})\)
and fix an order on \(V\), with \(V_{<v}\) its initial segment before
\(v\). Diminishing returns gives
\[
 \sum_{v\in P}a_v-h(P)\le\sum_{v\in P}
                   \bigl(a_v-h(v\mid V_{<v})\bigr).
\]
On weighting and summing, each nonnegative term on the right has
coefficient at most \(r\). The chain rule on \(V\) gives the claim.
The formal summation in \lean{TapeStorageRates.patch_deficit_packing}
uses the weighted submodular cover inequality from the Lean library
of~\cite{PaperTwo}.
\end{outline}

\begin{corollary}[Errors of local predictors]\label{cor:patch-prediction}
\formal{TapeStorageRates/PatchPrediction.lean}{TapeStorageRates.PatchPredictor.prediction_capacity,TapeStorageRates.PatchPredictor.prediction_deficit}
Let \(X\) be a random binary word on \(n=\card V\) coordinates.
Use unit patch weights and the ambient distance test above. Let \(r\)
bound the number of patches containing any coordinate. Each occurrence
\(i=(j,q)\) uses a fixed Boolean predictor of \(X_q\) from \(X_{B_{jq}}\)
with error probability \(e_i\). Assume the number \(m\) of scored occurrences is positive. Set
\(\bar e=m^{-1}\sum_i e_i\), and \(\delta=n-H(X_V)\).
With \(h_2(t)=-t\log_2t-(1-t)\log_2(1-t)\) and \(0\log_2 0=0\),
\[
 m[1-h_2(\bar e)]
 \le \sum_i[1-h_2(e_i)]
 \le r\bigl(\delta+\mathrm{TC}(X)\bigr)
 \le 2r\delta .
\]
\end{corollary}
\begin{outline}
The predictor and its error bit together determine \(X_q\). Thus
\(H(X_q\mid X_{B_{jq}})\le h_2(e_i)\), so \(c_q-h_2(e_i)\)
is a lower bound on the contribution of occurrence \(i\) in
\cref{thm:patch-credit} for every choice of the capacities. Use \(c_v=1\) in
\cref{thm:patch-credit} and \(H(X_v)\le1\). Concavity of binary
entropy gives the bound in terms of \(\bar e\).
\end{outline}

Equality holds at every error level \(0\le e\le1\).
In one patch, score both coordinates of \((U,U\mathbin{\oplus}Z)\),
where \(U\) is a fair bit and \(Z\) is an independent bit with
\(\Pr(Z=1)=e\). Each predictor copies its partner. Both paired
distances are infinite, \(r=1\), both errors equal \(e\), and the joint
entropy is \(1+h_2(e)\). Thus the two sides of the bound
\(m[1-h_2(\bar e)]\le2r\delta\) equal \(2[1-h_2(e)]\).
For \(e\ne 1/2\) this proves that the factor two is necessary for general binary laws. The example is \lean{TapeStorageRates.noisy_pair_prediction_sharp}.

For a uniform word in a binary code of length \(n>0\), suppose
\(\delta/n\le D/s\) with \(D,s>0\). If \(m\ge\alpha n\), then
\[
 r\ \ge\ \frac{\alpha[1-h_2(\bar e)]}{2D}\,s.
\]
This is \lean{TapeStorageRates.PatchPredictor.overlap_of_near_unit_rate}.
Predicting a positive density of separated targets with a fixed advantage
over a fair guess therefore requires overlap growing at least linearly
in \(s\). For the families constructed in \cref{thm:main}, take the radius
\(2w-1\) after fixing \(s\) and \(w\). For a native fibre supplied by
\cref{thm:native-tail}, \(X\) can be taken uniform on that particular
fibre, using its own actual permitted helpers and entropy deficit.

\begin{corollary}[Exact recovery over a finite alphabet]
\label{prop:patch-coverage}
\formal{TapeStorageRates/Patches.lean}{TapeStorageRates.recovery_patch_coverage}
Let \(C\subseteq\Sigma^V\) be nonempty, where \(a=\card{\Sigma}\ge2\).
Each patch \(P_j\) serves at most \(w\ge1\) targets in \(A_j\subseteq P_j\).
For each \(q\in A_j\), its value on \(C\) is determined by the
coordinates in \(N(q)\cap P_j\). Each coordinate belongs to at most
\(r\) patches. Let \(S\subseteq\bigcup_j A_j\), and let \(b_N(t)\)
count all coordinates in \(V\) at paired distance at most \(t\).
Then
\[
 \card S\le b_N(2w-1)+2r\bigl(n-\log_a\card C\bigr).
\]
In fact the last term alone bounds the number of distinct targets in
\(S\) at distance greater than \(2w-1\).
\end{corollary}
\begin{outline}
Give \(C\) the uniform law. In each patch score only the served targets
at distance greater than \(2w-1\). Their conditional entropies given the
permitted helpers vanish. In \cref{thm:patch-credit} take
\(c_v=\log_2 a\). Each long occurrence contributes \(\log_2 a\), and
every distinct long target in \(S\) contributes at least once. Both
global deficits are at most \(n\log_2 a-\log_2\card C\).
\end{outline}

Thus even if the alphabet varies with the code, a family with a
vanishing fraction of short pairs at radius \(2w-1\) has upper
limiting covered fraction at most \(2r(1-\liminf_i R_i)\), where
\(R_i=\log_{\card{\Sigma_i}}\card{C_i}/\card{V_i}\) and \(w,r\)
are fixed. This is \lean{TapeStorageRates.patch_coverage_limsup}.
Disjoint repeated pairs attain the finite bound at \(r=1\), as in
\lean{TapeStorageRates.repetition_pairs_sharp}.

\begin{corollary}[Input-dependent choice of predictors]\label{cor:patch-menu}
\formal{TapeStorageRates/PatchPredictionEvent.lean}{TapeStorageRates.patch_prediction_menu_probability,TapeStorageRates.patch_prediction_event_probability}
Let \(E\subseteq\{0,1\}^V\), \(n=\card V\), \(0\le e\le1/2\), and
\(r\ge1\). Fix a menu of at most \(2^h\) predictor systems satisfying
the ambient distance test with common bounds \(w,r\). A system specifies
its relation, patches, scored targets, helpers and prediction rules.
Suppose that, for each \(x\in E\), some system scores at least \(m\)
occurrences and errs on at most a fraction \(e\) of all the occurrences
it scores on \(x\). Then
\[
 \frac{\card E}{2^n}\le
 2^{\,h-m[1-h_2(e)]/(2r)}.
\]
In particular a transcript selecting the system can be charged
\(h=\log_2\card{\{\text{labels attained on }E\}}\) when \(E\ne\varnothing\).
\end{corollary}
\begin{outline}
For each system, restrict to the inputs of \(E\) on which it meets
the two requirements. Under the uniform law on any such nonempty subset,
its mean error is at most \(e\). By \cref{cor:patch-prediction} and
monotonicity of \(h_2\) on \([0,1/2]\), the subset has at most
\(2^{\,n-m[1-h_2(e)]/(2r)}\) elements. Their union contains \(E\).
\end{outline}

A tape lower bound along this route still requires a proof that
successful executions admit predictors with the stated accuracy and
overlap. On a fixed native fibre one can test the predictors under its
uniform law. On the full cube one must also bound the number of distinct
systems selected by the trace, the genie word and any other observation
used to choose the targets.

\section{What remains open}\label{sec:open}

\paragraph{How close to one can the abstract rate get?}
Known binary constructions exceed $0.5013$ at a fixed degree cap.
For larger caps the large-prime construction gives guaranteed rates
arbitrarily close to $4-2\sqrt3$ from below
\cite[Theorem~A$^*$, Lemma~8.1 and Corollary~A$'$]{PaperThirteen}.
Both facts leave room for the rate upper bound in \cref{cor:gap}.
A proof along this route must control all admissible families at the
caps used in the test. A concrete construction target for a given tape
count $k$ is a family of limiting lower rate at least
$1-K f_{4k^2+1}(\Delta)$ at each sufficiently large cap, with one
constant $K$. Such families would preclude the gap in \cref{cor:gap}
for that $k$. A tape algorithm would also have to realize their
recoveries by one causal computation within its running-time budget.

\paragraph{Which helpers can a physical computation use?}
\Cref{thm:native-tail} keeps the time at which each retained bit is
appended. This allows a prospective impossibility proof to choose
patches from a tape schedule and test them on its actual code. Fix a positive prediction advantage $\xi$ and a target
density $\alpha>0$. For each simulator batch size $s$, choose a per-patch
target bound $w_s$ before the dimension grows and use paired-distance
radius $2w_s-1$. Consider cells with the following property. Under the
uniform law on that cell, one fixed system scores at least an $\alpha$
fraction of its retained coordinates and predicts them from permitted
helpers with mean accuracy at least $1/2+\xi$. Its overlap is bounded by
a common $r_s=o(s)$. Must such cells carry a fraction of the uniform input
mass bounded away from zero as the dimension grows?

For a full successful cube at fixed $s$, the cells of deficit fraction
$O(1/s)$ carry all but at most $2^{-L/s}$ of the input mass. A positive
mass bound independent of the dimension would give a cell meeting both
conditions. On that cell \cref{cor:patch-prediction} would require overlap
of order $s$, contradicting the proposed upper bound at large $s$. If
targets or predictors are chosen separately for each input, the menu
test of \cref{cor:patch-menu} charges the number of systems used.
Establishing a prediction advantage on actual runs remains open.
Two independent fair helpers can determine their XOR exactly even though
either one alone gives chance accuracy.

\appendix
\section{Formal statements and verification}\label{sec:verification}

The accompanying Lean development is authoritative for the numbered
statements and their proofs. The human-readable outlines give the
mathematical arguments and explain the hypotheses. Inputs, moves, output
and halting use the first paper's frozen machine model \cite{PaperOne}.

\paragraph{Packet coordinates.}
\lean{TapeStorageRates.PaidInputs.GrowingWidthSuccess} quantifies over a common
choice of dimension and entry width for each size threshold and accuracy
tolerance. The source machine supplies its own class of successful inputs.
\lean{TapeStorageRates.PaidClasses.compiler} realizes the eleven-tick
batch computation. \lean{TapeStorageRates.packet_code_at_threshold}
performs the finite extraction and
\lean{TapeStorageRates.family_of_packet_classes_at_threshold} chooses
dimensions to give a family. The positive lengths, two-sided degree
bounds and paired distances belong to
\lean{TapeStorageRates.SeparatedBinaryFamily}. The comparison with the
supremum of limiting upper rates is
\lean{TapeStorageRates.RateWitnesses.extremal_rates}.

\paragraph{Native coordinates.}
\lean{TapeStorageRates.NativeTypicalPartition} chooses one cut and
retained coordinate set per attained trace. Each input restricts to a
word of its own cell code. The membership proof is
\lean{TapeStorageRates.NativeTypicalPartition.projected_mem}. The tail theorem concerns the uniform
probability on $G$ under these assignments. For \cref{thm:native-tail},
depth reduction is applied to the coarse work walks; the joint support
of complete batch records pays for the physical revelations.
\Cref{cor:native} instead uses the native coding hypothesis of the
fourteenth paper at its specified physical step width and depth
\cite[Corollary~L]{PaperFourteen}.

\paragraph{Predictions.}
\lean{TapeStorageRates.PatchPredictor.missed} masks all inputs outside
a predictor's permitted helpers before evaluating its guess.
The development uses natural logarithms for Shannon entropy. Dividing
by $\ln 2$ gives the bit units of \cref{sec:patches}. Counting one scored
occurrence for each patch and target allows one target to occur in
several patches. The finite-alphabet exact-coverage result counts
distinct covered targets.

\begin{corollary}[The original qualitative target]\label{cor:formal}
\formal{TapeStorageRates/Main.lean}{TapeStorageRates.qualitative_transfer}
Under \textnormal{(T)}, for each $\varepsilon>0$ and every sufficiently large
cap there is an admissible family of limiting lower rate at least
$1-\varepsilon$.
\end{corollary}
\begin{outline}
A witness to (T) succeeds at all canonical dimensions. The original target uses $\max(1,k)$ in place of the actual tape count
$k$. The quantitative statement is proved by \lean{TapeStorageRates.quantitative_transfer}. Its qualitative consequence follows by \lean{TapeStorageRates.qualitative_of_quantitative}. At that fixed tape
count the deficit scale tends to zero. \Cref{thm:main} also gives the
bound with $k$ itself, including $k=0$.
\end{outline}

\paragraph{Reproducing the checks.}
Use the complete repository so the sibling Lean packages are available.
The project pins Lean and Mathlib v4.34.1. From this paper's directory run
\begin{verbatim}
(cd lean && lake build && lake env lean PaperAxioms.lean)
python3 tools/lean_index.py --check
\end{verbatim}
The generated axiom file queries the declarations cited in the text and
in the formal status lines. Their dependencies use at most
\texttt{propext}, \texttt{Classical.choice} and \texttt{Quot.sound}.
The finite phase calculations use locally increased heartbeat limits;
their proof terms are checked by the same kernel.
The README gives the preparation steps for a fresh checkout.

% Generated by tools/lean_index.py from the status lines; do not edit by hand.
\section{Index of Lean declarations}\label{sec:leanindex}

The status lines of the paper cite the following declarations, listed by Lean module
with their fully qualified names. Module files are relative to the paper's
\path{lean/} directory. The axiom checks in \path{lean/PaperAxioms.lean} are generated
from the same sources.

\paragraph{\texttt{TapeStorageRates.Main}}
{\small\raggedright\leanlist{TapeStorageRates.exponentially_small_at_widths_of_sparse_gap,TapeStorageRates.qualitative_transfer,TapeStorageRates.rate_witnesses_of_full,TapeStorageRates.rate_witnesses_of_growing_widths,TapeStorageRates.rate_witnesses_of_near_full}.\par}

\paragraph{\texttt{TapeStorageRates.NativeCompiler}}
{\small\raggedright\leanlist{TapeStorageRates.Native.same_partial_function,TapeStorageRates.Native.source_support,TapeStorageRates.Native.source_support_fixed}.\par}

\paragraph{\texttt{TapeStorageRates.NativeTypical}}
{\small\raggedright\leanlist{TapeStorageRates.NativeTypicalPartition.deficit_moment,TapeStorageRates.NativeTypicalPartition.tail_of_joint_count}.\par}

\paragraph{\texttt{TapeStorageRates.NativeTypicalRates}}
{\small\raggedright\leanlist{TapeStorageRates.typical_native_rate_tail,TapeStorageRates.typical_native_rates_full}.\par}

\paragraph{\texttt{TapeStorageRates.PacketFamily}}
{\small\raggedright\leanlist{TapeStorageRates.packet_code_at_threshold}.\par}

\paragraph{\texttt{TapeStorageRates.PatchPrediction}}
{\small\raggedright\leanlist{TapeStorageRates.PatchPredictor.prediction_capacity,TapeStorageRates.PatchPredictor.prediction_deficit}.\par}

\paragraph{\texttt{TapeStorageRates.PatchPredictionEvent}}
{\small\raggedright\leanlist{TapeStorageRates.patch_prediction_event_probability,TapeStorageRates.patch_prediction_menu_probability}.\par}

\paragraph{\texttt{TapeStorageRates.PatchResidual}}
{\small\raggedright\leanlist{TapeStorageRates.patch_capacity_of_long_targets,TapeStorageRates.weighted_patch_capacity}.\par}

\paragraph{\texttt{TapeStorageRates.Patches}}
{\small\raggedright\leanlist{TapeStorageRates.recovery_patch_coverage}.\par}

\paragraph{\texttt{TapeStorageRates.TraceTransfer}}
{\small\raggedright\leanlist{TapeStorageRates.native_frozen_trace_exclusion,TapeStorageRates.native_tests_exclude_linear}.\par}

\paragraph{Also named in the text}
{\small\raggedright From this paper's library: \leanlist{TapeStorageRates.BatchCoarsening.Coarsening.path_length_le,TapeStorageRates.Native.control_bits,TapeStorageRates.Native.copy_trace_card,TapeStorageRates.Native.simulates_halt,TapeStorageRates.Native.trace_support_rate_limsup,TapeStorageRates.NativeBatchRecord.compiler_joint_record_support,TapeStorageRates.NativeFibreWitness.outputTime_eq_appendTime,TapeStorageRates.NativeTypicalCut.mem_code_iff,TapeStorageRates.NativeTypicalCut.storage,TapeStorageRates.NativeTypicalPartition,TapeStorageRates.NativeTypicalPartition.projected_mem,TapeStorageRates.PaidClasses.compiler,TapeStorageRates.PaidInputs.GrowingWidthSuccess,TapeStorageRates.PatchPredictor.missed,TapeStorageRates.PatchPredictor.overlap_of_near_unit_rate,TapeStorageRates.Prefix.control_bits,TapeStorageRates.Prefix.simulates_halt,TapeStorageRates.Prefix.source_support,TapeStorageRates.Prefix.time_le_eleven,TapeStorageRates.RateWitnesses.extremal_rates,TapeStorageRates.SeparatedBinaryFamily,TapeStorageRates.aligned_low_trace_dimensions,TapeStorageRates.elementary_gap_ratio_tendsto_zero,TapeStorageRates.family_of_packet_classes_at_threshold,TapeStorageRates.genie_scaled_cap_degree_bound_at,TapeStorageRates.native_batch_size_for_cap,TapeStorageRates.native_five_eighths,TapeStorageRates.native_joint_typical_partition,TapeStorageRates.noisy_pair_prediction_sharp,TapeStorageRates.patch_coverage_limsup,TapeStorageRates.patch_deficit_packing,TapeStorageRates.qualitative_of_quantitative,TapeStorageRates.quantitative_transfer,TapeStorageRates.ratio_tendsto_of_power_gap,TapeStorageRates.repetition_pairs_sharp}.\par}
{\small\raggedright From Paper 11's library: \leanlist{MergeDepth.skel_paired_packing}.\par}

\begingroup
\raggedright
\bibliographystyle{plainurl}
\bibliography{references}
\endgroup
\end{document}
